\par\medskip\noindent
\textbf{\SGAInlineRefTarget{sga3:xix:paragraph:4-5}4.5.}---
Let \(S\) be a scheme, let \(T\) be an \(S\)-torus, let \(P\) be an
\(S\)-group with operator group \(\mathbb G_{m,S}\)
\textup{(for example, a vector bundle)}, and let \(\alpha\) be a
character of \(T\).  We write \(T\mathbin{\cdot_\alpha}P\) for the
semidirect product of \(P\) by \(T\), where \(T\) acts on \(P\)
through the composite morphism
\[
  T
  \xrightarrow{\ \alpha\ }
  \mathbb G_{m,S}
  \longrightarrow
  \underline{\operatorname{Aut}}_{S\text{-gr.}}(P).
\]
\SGARefTarget{sga3:xix:paragraph:4:5}

\par\medskip\noindent
\SGAInlineRefTarget{sga3:xix:definition:4-6}\textbf{\textit{\SGARefLink{sga3:xix:definition:4:6}{Definition 4.6}.}}---
\textit{Let \(S\) be a scheme, let \(G\) be an \(S\)-group scheme,
let \(T\) be a subgroup of \(G\), let \(\alpha\) be a character of
\(T\), and let \(L\) be an \(\mathcal O_S\)-module.  Let}
\[
  p\colon\mathbb W(L)\longrightarrow G
\]
\textit{be a morphism of \(S\)-group functors.}\footnote{Editor's
note: the edition uses a bold \(\mathbf W(L)\) here because, for an
arbitrary \(\mathcal O_S\)-module \(L\), the functor need not be
representable.  Below, however, \(L\) is assumed finite locally free
\textup{(indeed, invertible)}, in which case the functor is
represented by the vector bundle \(\mathbb V(L^\vee)\), and the
notation \(\mathbb W(L)\) is used without this distinction.}
\textit{We say that \(p\) is \emph{normalized by \(T\) with
multiplier \(\alpha\)} if it satisfies the following equivalent
conditions:}

\begin{enumerate}
\item[\textup{(i)}]
\(p\) is a morphism of objects with operator group \(T\), where \(T\)
acts on \(\mathbb W(L)\) through \(\alpha\), and on \(G\) by inner
automorphisms.  In other words, for every \(S'\to S\), every
\(t\in T(S')\), and every
\[
  x\in\mathbb W(L)(S')
  =
  \Gamma(S',L\otimes_{\mathcal O_S}\mathcal O_{S'}),
\]
one has
\[
  \operatorname{int}(t)p(x)=p\bigl(\alpha(t)x\bigr).
\]

\item[\textup{(ii)}]
The morphism
\[
  T\mathbin{\cdot_\alpha}\mathbb W(L)\longrightarrow G,
  \qquad
  (t,x)\longmapsto t\,p(x),
\]
defined by multiplication in \(G\), is a group morphism.
\end{enumerate}
\SGARefTarget{sga3:xix:definition:4:6}

\par\medskip\noindent
\SGAInlineRefTarget{sga3:xix:lemma:4-7}\textbf{\textit{\SGARefLink{sga3:xix:lemma:4:7}{Lemma 4.7}.}}---
\textit{Under the hypotheses of \SGARefLink{sga3:xix:definition:4-6}{4.6}, suppose in addition that \(G\)
is smooth with connected fibers,\footnote{Editor's note: compare
editor's note \textup{(28)}.} that \(T\) is a maximal torus of \(G\),
and that \(L\) is invertible.  If \(p\) is a monomorphism and
\(\alpha\) is nonzero on every fiber, then \(\alpha\) is a root of
\(G\) with respect to \(T\).}
\SGARefTarget{sga3:xix:lemma:4:7}

Indeed,
\[
  \operatorname{Lie}(p)\colon L\longrightarrow\mathfrak g
\]
is a monomorphism which factors through \(\mathfrak g^\alpha\).

\par\medskip\noindent
\SGAInlineRefTarget{sga3:xix:proposition:4-8}\textbf{\textit{\SGARefLink{sga3:xix:proposition:4:8}{Proposition 4.8}.}}---
\textit{Under the hypotheses of \SGARefLink{sga3:xix:lemma:4-7}{4.7}, suppose that \(G\) is reductive
and that \(p\) is a monomorphism.  Then \(\alpha\) is a root of \(G\)
with respect to \(T\), and \(\operatorname{Lie}(p)\) induces an
isomorphism}
\[
  \operatorname{Lie}(p)\colon
  L\xrightarrow{\ \sim\ }\mathfrak g^\alpha.
\]
\SGARefTarget{sga3:xix:proposition:4:8}

Indeed, by \SGARefLink{sga3:xix:lemma:4-7}{4.7} and the fact that \(\mathfrak g^\alpha\) is
invertible, it is enough to prove that \(\alpha\) is nonzero on every
fiber.  Let \(s\in S\) be such that
\(\alpha_{\bar s}=0\) \textup{(\(=1\) in multiplicative notation)}.
If \(X\) is a nonzero section of
\(L\otimes_S\kappa(\bar s)\), then \(p(X)\) is a unipotent element
different from the identity in \(G(\bar s)\) which centralizes
\(T_{\bar s}\).  This is impossible, since \(T_{\bar s}\) is its own
centralizer.

\par\medskip\noindent
\SGAInlineRefTarget{sga3:xix:corollary:4-9}\textbf{\textit{\SGARefLink{sga3:xix:corollary:4:9}{Corollary 4.9}.}}---
\textit{Under the hypotheses of \SGARefLink{sga3:xix:proposition:4-8}{4.8}, there exists a monomorphism of
groups with operator group \(T\) \textup{(that is, normalized by
\(T\) with multiplier \(\alpha\))}
\[
  \mathbb W(\mathfrak g^\alpha)\longrightarrow G
\]
which induces on Lie algebras the canonical morphism
\(\mathfrak g^\alpha\to\mathfrak g\).}
\SGARefTarget{sga3:xix:corollary:4:9}

We shall soon see that \SGARefLink{sga3:xix:corollary:4-9}{4.9} in fact holds whenever \(G\) is a
reductive group and \(\alpha\) is a root of \(G\) with respect to
\(T\), and that such a morphism is unique.

